% !PRINTED 131
\weberpage{131}
If we can now also prove that there are in $U$ \gesperrt{exactly $(p-1)\nu$ elements whose degree is divisible by $p$,} then we have reached our goal; for it then follows that among the elements $U$, and hence also among the elements of $P$, there are exactly
\[
p\nu - (p-1)\nu = \nu
\]
whose degree is a divisor of $\nu$.

But this results from the following:

$\beta$) \gesperrt{If $v$ is an element of $V$, then all the elements $vQ$ (whose number is $p$) belong to $V$. Among them there are $p-1$ whose degree is divisible by $p$, and one whose degree is not divisible by $p$. In $V$ there are $(p-1)r_2$ elements whose degree is divisible by $p$, and $r_2$ elements whose degree is not divisible by $p$.}

For if $a$ is an element of $Q$ different from $1$, then, since $v$ belongs to $R$, so that $v^{-1}av$ is contained in $Q$, and since $Q$ consists of the powers of $a$,
\weq{1}{v^{-1}av = a^s.}
By repeated application there results from this, for every exponent $h$,
\weq{2}{v^{-h}av^h = a^{s^h}.}
If we take for $h$ the degree of the element $v$, which by the hypothesis is a divisor of $p\nu$, and thus set $v^h = 1$, it follows that
\weq{3}{a = a^{s^h}, \quad s^h \equiv 1 \ (\mathrm{mod}\ p).}
But no prime factors divide $h$ other than such as are equal to or greater than $p$, and hence $p-1$ is relatively prime to $h$, and one can satisfy the congruence
\weq{4}{xh \equiv 1 \ (\mathrm{mod}\ p-1)}
accordingly, from the \gesperrt{Fermat's} theorem there results:
\weq{5}{s^{xh} \equiv s \equiv 1 \ (\mathrm{mod}\ p),}
hence by (1):
\weq{6}{av = va,}
i.e.\ $a$ and $v$ are permutable. From this it follows, for every exponent $\lambda$,
\weq{7}{(va)^\lambda = v^\lambda a^\lambda.}
